\documentclass[aspectratio=169]{beamer}
\usepackage[T1]{fontenc}
\usepackage{graphicx}
\usepackage{xcolor}
\graphicspath{{figures/common/}{figures/01/}}
\definecolor{vaftblue}{RGB}{25,76,127}
\setbeamercolor{structure}{fg=vaftblue}
\setbeamerfont{title}{size=\fontsize{50}{56}\selectfont}
\setbeamerfont{subtitle}{size=\fontsize{24}{28}\selectfont}
\setbeamerfont{frametitle}{size=\fontsize{35}{40}\selectfont}
\setbeamerfont{normal text}{size=\fontsize{18}{22}\selectfont}
\setbeamertemplate{navigation symbols}{}
\setbeamertemplate{footline}[frame number]
\title{Getting Started with VAFT}
\subtitle{VAFT Introductory Tutorial | Session 01}
\author{}
\date{}

\begin{document}
\begin{frame}
  \titlepage
\end{frame}

\begin{frame}{Why this session exists}
  \begin{itemize}
    \item Read a tokamak shot as a time-resolved experiment.
    \item Find diagnostic data in an OMAS/IMAS-style IDS structure.
    \item Make and question a plot through VAFT's public API.
  \end{itemize}
\end{frame}

\begin{frame}{A shot is a time-resolved experiment}
  \begin{enumerate}
    \item Prefill and coil preparation establish initial conditions.
    \item Breakdown and current rise mark the start of plasma formation.
    \item A sustained-current interval supports diagnostic interpretation.
    \item Ramp-down and shutdown complete the discharge.
  \end{enumerate}
  \vspace{0.5em}
  A trace suggests questions about these phases; it is not a complete physical
  explanation on its own.
\end{frame}

\begin{frame}{Diagnostics give complementary views}
  \begin{itemize}
    \item Magnetic measurements track plasma current, fields, and flux.
    \item PF-coil currents document the programmed actuator response.
    \item UV spectroscopy adds a time-resolved impurity-emission signal.
  \end{itemize}
  \vspace{0.6em}
  Compare timing first. Use channel meaning, calibration, and other diagnostics
  before making a causal claim.
\end{frame}

\begin{frame}{From IDS to a plot}
  \centering
  \Large
  VEST shot \(\rightarrow\) OMAS ODS \(\rightarrow\) IDS root \(\rightarrow\)
  documented path \(\rightarrow\) VAFT public plot recipe
  \vspace{0.7em}

  \normalsize
  Examples: \texttt{magnetics}, \texttt{pf\_active},
  \texttt{spectrometer\_uv}, and \texttt{equilibrium}.
\end{frame}

\begin{frame}{The reproducible offline route}
  \begin{enumerate}
    \item Load VAFT's packaged sample ODS for shot 39915.
    \item Inspect IDS roots and the \texttt{magnetics.time} and
      \texttt{magnetics.ip.0.data} paths.
    \item Discover recipes with \texttt{vaft.omas.available\_plots}.
    \item Plot plasma current, a magnetic overview, and UV intensity.
  \end{enumerate}
\end{frame}

\begin{frame}{The guided lab route}
  Set \texttt{VAFT\_TUTORIAL\_MODE=lab} after configuring HSDS with
  \texttt{hsconfigure}. The notebook then performs a read-only magnetic-data
  selection with \texttt{vaft.database.open}.
  \vspace{0.7em}

  No credentials or network are required for the offline lesson. An unavailable
  lab connection is a clear skipped extension, not an error in the core work.
\end{frame}

\begin{frame}{Ask a defensible first question}
  \begin{itemize}
    \item Where does plasma current rise, persist, and return toward zero?
    \item Which magnetic or spectrometer feature occurs near that transition?
    \item What metadata or additional diagnostic would strengthen the interpretation?
  \end{itemize}
\end{frame}

\begin{frame}{Independent exercise}
  Choose a supported plot recipe other than the guided plasma-current trace.
  Resolve its exported \texttt{vaft.omas.plot\_*} function, then record:
  \begin{enumerate}
    \item the diagnostic and time interval;
    \item one observed feature; and
    \item one missing piece of evidence needed for interpretation.
  \end{enumerate}
\end{frame}

\begin{frame}{Next: operation and vacuum fields}
  Session 02 connects these first diagnostic traces to the programmed discharge
  scenario, PF-coil evolution, and vacuum-field interpretation.
\end{frame}
\end{document}
